\section{Empirical detail}\label{app:emp}

\subsection{Data acquisition and panel construction}

The Federal Audit Clearinghouse publishes the complete structured results of
every Single Audit submission as bulk files at
\url{https://www.fac.gov/data/download/}, without registration,
authentication or fee. Seven files were retrieved in September 2026 and cover
fiscal years 2016 onward. No free-text finding narratives, no restricted data
and no nonpublic records of any kind were used. The regulatory provisions
quoted are from Title~2 of the Code of Federal Regulations; the 2023 annual
edition is the one in force over the sample period.

The unit of analysis is the auditee--program--year. A ``program'' is the
regulatory program unit: a cluster of programs where one exists, otherwise a
single assistance listing number. That follows 2~CFR 200.518(b)(3), under
which ``a cluster of programs is treated as one program'' for the Type~A
determination, so the unit matches the one the rules operate on. The window
is fiscal years 2016 through 2023, restricted to audits whose fiscal year began
within the period governed by the \$750,000 audit and Type~A thresholds in
200.501(a) and 200.518(b)(1) as they stood before the 2024 revision, which
raised both to \$1,000,000 for fiscal years beginning on or after 1~October
2024 \citep{omb2024}. The panel ends before that revision takes effect, so a
single threshold regime governs throughout.

Table~\ref{tab:sample} shows how the panel narrows and why a missing prior
year is not coded as a clean one. The panel distinguishes a resolved
prior-year observation with no finding from the absence of any prior-year
record for that program: of the \nPanelRows{} program-years, \nSone{} have a
resolved prior-year finding and \nStwo{} a resolved prior-year clean result;
the remainder are excluded rather than recoded. Sample~A restricts further to
program-years in which the program was audited as major in both years. It
holds the major-program-selection channel fixed. It does not hold audit
intensity or condition-specific observation fixed, and cannot, for the reason
given in Section~\ref{sec:regime}: the follow-up requirement applies within
Sample~A.

Two definitions of ``repeat'' are used. The primary one is the auditor's own
coding. The secondary one is a validated historical link: a finding whose
cited prior reference number appears as a finding reference recorded for the
same auditee and program in the immediately preceding year. Of the
\nLinkFlagged{} findings that cite a prior-finding reference, \nLinkResolved{}
resolve this way, \nLinkPct{}~percent. Non-resolution has innocent causes: the
prior finding may sit two or more years back, may be recorded against a
different program unit, or may carry a renumbered reference. Neither measure
is a gold standard, so both are reported. Intervals throughout are 95 percent
cluster bootstrap with \nBootReps{} replications, resampling \nNAuditees{}
auditees together with their full histories, seed \nBootSeed{}.

\subsection{Every specification}

Table~\ref{tab:specs} collects the risk differences under all five
specifications reported anywhere in this paper. The first row is the primary
result of Table~\ref{tab:fourstate}. The second substitutes the validated
historical link for the auditor's repeat flag. This definition mechanically
forces both repeat states to zero for program-years with no prior finding, so
its risk differences are interpretable and its risk ratios are not. It has a
second consequence, which the table makes explicit. Under the auditor's flag,
every recorded finding is either marked a repeat or is not, so the four
outcome states exhaust the possibilities and their risk differences add to the
difference in ``any finding''. Under the validated link they do not. A
program-year whose findings were all flagged as repeats, none of which
resolves to a same-program prior-year reference, has recorded no new finding
and no \emph{validated} repeat, and therefore falls outside all four states.
That fifth state occurs in \nVlinkFifthNoPrior{} of program-years with no
prior finding and \nVlinkFifthPrior{} of those with one, a difference of
\nVlinkFifth{}, which is exactly the gap between the four state columns and
``any finding'' in that row. It is given its own column rather than absorbed
anywhere, and the sum is checked in every row.

The third row drops the requirement that the program be audited as major in
both years, using all program-years with a resolved prior observation. The
fourth restricts to transitions whose Clearinghouse provenance is stable across
the two years, which addresses the possibility that repeat-flag capture differs
between the two administering systems that operated during the window; that
restriction excludes \nProvExcluded{} transitions. The fifth uses the
three-year cumulative history rule, requiring resolved observations at $t-1$,
$t-2$ and $t-3$. Across all five, the any-new margin is the most stable
quantity in the table.

\subsection{The allocation of mixed periods}

Table~\ref{tab:zeta} gives the full sensitivity of the ``percent due to
recurrence'' statistic to the treatment of program-periods that recorded both a
repeat and a new finding. $\zeta$ is the fraction of mixed periods counted as
recurrence. The $\zeta = 1$ row gives \nNineZeroNine{}~percent; the $\zeta = 0$
row gives \nZetaLo{}~percent. The arithmetic is not in doubt: the numerator
$\nDRonly{} + \zeta \cdot \nDBoth{}$ and the denominator \nDAny{} are read
directly off Table~\ref{tab:fourstate}. A correctly computed statistic can
still be uninformative, because the quantity it reports is fixed by the
allocation rule rather than by anything observed.

\subsection{Compliance-requirement categories}

The Clearinghouse records a requirement category for every reported
deficiency. It does not record which requirements applied to, or were tested
in, a given program-period; that universe is defined by the Compliance
Supplement, which is external to these data. A first-in-area hazard therefore
has no constructible denominator here, and none was constructed.

What can be said is descriptive and applies only to programs that recorded
findings in both years. Among those \nOverlapN{} Sample~A program-years,
\nOverlapShares{}~percent share at least one requirement category across the
two years; the mean Jaccard overlap between the two category sets is
\nOverlapJaccard{}; on average \nOverlapEntering{} categories enter and
\nOverlapExiting{} exit between the two years. These are descriptions of a
conditioned subset, not comparisons across history groups, and they carry no
implication about the rate at which problems appear in new requirement areas.

\subsection{Horizons beyond one year}

Results at $t+2$ and $t+3$ are not reported in the main text, and the reason
is structural rather than statistical. Following a program past $t+1$ requires
it to be audited as a major program again, and that is exactly the quantity
the regulation makes depend on the finding record: a prior finding can block
low-risk Type~A status and so make recoverage more likely. The programs that
can be followed are therefore not a random subset of the programs one would
want to follow, and the selection runs in the direction that would flatter a
persistence interpretation.

The attrition is severe. Of \nHorizonNone{} Sample~A finding-years with
resolved follow-up at $t+1$, \nHorizonNtwo{} can be followed to $t+2$ and
\nHorizonNthree{} to $t+3$, losses of \nHorizonLossTwo{} and
\nHorizonLossThree{}~percent. Across those shrinking and increasingly selected
samples the explicit-repeat rate falls from \nHorizonRepeatOne{} at $t+1$ to
\nHorizonRepeatThree{} at $t+3$. That decline is what both a decay story and a
selection story predict, and the design cannot separate them, which is why the
numbers are recorded here rather than interpreted.

\subsection{Data and code availability}

Reproduction requires the seven bulk files, a build of the
auditee--program--year panel following the definitions above, and estimation
of the outcome shares with cluster-bootstrap intervals as described. The panel
construction, the four-state decomposition, the cluster bootstrap and the
simulations reported in Sections~\ref{sec:design} and \ref{sec:scoring} are
implemented in Python and DuckDB with a fixed master seed. The quantities
reported from simulation are either exact calculations or Monte Carlo results
from a single prespecified configuration; no parameter search was performed,
and each table states which of the two it is. The full analysis code, the
intermediate panel and cryptographic digests of every input and output will be
deposited with the journal on conditional acceptance and retained for at least
seven years.

\begin{table}[tbp]
\centering
\caption{\textbf{How the sample was built, and what was left out.} The upper
panel narrows the panel to program-years in which the same program can be
observed in two consecutive years with its major-program status known in both.
The lower panel is the reason the paper does not simply code a missing prior
year as a clean year. In 366,506 program-years --- 15.35 percent of the panel --- the
auditee filed an audit last year but this program does not appear in it. Those
program-years record a finding 3.26 percent of the time, which is not the rate of
a clean program. They are excluded rather than recoded, because the absence of a
record is not a clean result.}
\label{tab:sample}
\footnotesize
\setlength{\tabcolsep}{3pt}
\begin{tabular}{lrr}
\toprule
Step & Program-years & Excluded here \\
\midrule
All program-years in the panel & 2,387,511 &  \\
Within the \$750,000 regime (FY2016--2023) & 2,387,511 & 0 \\
With a resolved prior-year observation & 1,600,898 & 786,613 \\
Sample A: audited as a major program in both years & 206,926 & 1,393,972 \\
\bottomrule
\end{tabular}

\vspace{1.1em}

\footnotesize
\begin{tabular}{lrrrl}
\toprule
What the prior year looks like & Program-years & Share & \shortstack{Finding rate\\this year} & Treatment \\
\midrule
Prior year resolved, finding recorded & 66,131 & 2.77\% & 0.4723 & kept \\
Prior year resolved, no finding recorded & 1,534,767 & 64.28\% & 0.0206 & kept \\
Auditee filed last year; this program is not in that filing & 366,506 & 15.35\% & 0.0326 & excluded \\
No filing at all for this auditee last year & 160,861 & 6.74\% & 0.0706 & excluded \\
First year of the window (2016) & 259,246 & 10.86\% & 0.0462 & excluded \\
\bottomrule
\end{tabular}
\end{table}

\begin{table}[htbp]
\centering
\caption{Risk differences under every specification reported.}
\label{tab:specs}
\scriptsize
\setlength{\tabcolsep}{3pt}
\begin{tabular}{lrcccccc}
\toprule
& & \multicolumn{4}{c}{Mutually exclusive outcome states} & & \\
\cmidrule(lr){3-6}
Specification & $n$ & repeat only & new only & both & \shortstack{flagged repeat,\\ not linked} & any finding & any new finding \\
\midrule
Sample A, auditor flag (primary) & 206,926 & $+0.3073$ & $+0.0456$ & $+0.1488$ & --- & $+0.5017$ & $+0.1945$ \\
Sample A, validated link & 206,926 & $+0.2994$ & $+0.0491$ & $+0.1454$ & $+0.0079$ & $+0.5017$ & $+0.1945$ \\
All resolved transitions & 1,600,898 & $+0.2869$ & $+0.0617$ & $+0.1031$ & --- & $+0.4517$ & $+0.1648$ \\
Stable provenance & 167,017 & $+0.2935$ & $+0.0510$ & $+0.1439$ & --- & $+0.4885$ & $+0.1949$ \\
Three-year history & 94,360 & $+0.1959$ & $+0.0774$ & $+0.0942$ & --- & $+0.3674$ & $+0.1716$ \\
\bottomrule
\end{tabular}
\tablenotes{The primary row is
the one in Table~\ref{tab:fourstate}. The validated-link row mechanically forces
the repeat states to zero for program-years with no prior finding, so its risk
differences are usable and its ratios are not. It also opens a fifth outcome
state: a period whose findings were all auditor-flagged repeats, none of which
resolves to a same-program prior-year reference, records no new finding and no
validated repeat, and so falls outside all four. That column is empty for every
other specification, where the four states do exhaust the possibilities. In each
row the four state columns sum to ``any finding'', which the exhibit script
asserts on the unrounded values; the columns here are rounded independently, so
a row may add up one unit in the last place away from the total shown. The
any-new margin is the quantity that moves least across all five
specifications.}
\end{table}

\begin{table}[htbp]
\centering
\caption{The recurrence share as a function of a reporting convention. $\zeta$ is
the fraction of mixed program-periods --- those recording both a repeat and a
new finding --- counted as recurrence. Counting all of them gives
90.9~percent; counting none gives
61.2~percent. Nothing observable differs between the
two ends, which is why the paper reports the three-way split instead.}
\label{tab:zeta}
\small
\begin{tabular}{ccc}
\toprule
$\zeta$ & ``due to recurrence'' & ``first recorded'' \\
\midrule
0.00 & 61.2\% & 38.8\% \\
0.25 & 68.7\% & 31.3\% \\
0.50 & 76.1\% & 23.9\% \\
0.75 & 83.5\% & 16.5\% \\
1.00 & 90.9\% & 9.1\% \\
\bottomrule
\end{tabular}
\end{table}

